% =====================================================================
% === 1. IMPOSTAZIONI DI BASE (Codifica, Font, Lingua, Layout)
% =====================================================================
% NOTA: l'ordine inputenc → fontenc → lmodern → babel è obbligatorio.

\usepackage[utf8]{inputenc}   % Codifica sorgente UTF-8
\usepackage[T1]{fontenc}      % Codifica output T1 (necessario per accentate, copia da PDF, ricerca)
\usepackage{lmodern}          % Font Latin Modern in T1 (sostituisce i bitmap di default)
\usepackage[italian]{babel}   % Supporto lingua italiana (sillabazione, didascalie, ecc.)
\usepackage{microtype}  % aggiustamenti tipografici micro (elimina la maggior parte degli Overfull \hbox piccoli)

\usepackage{geometry}
\geometry{
    a4paper,
    margin     = 1in,
    headheight = 14pt   % fancyhdr con font 11pt richiede almeno 13.6pt; geometry compensa topmargin
}
\raggedbottom   % non forza l'allineamento del fondo pagina (consigliato per documenti tecnici)

% =====================================================================
% === 2. LAYOUT: TITOLI, INTESTAZIONE E INDICE
% =====================================================================

\usepackage{titlesec}         % Personalizzazione stile titoli di capitolo/sezione

%%%%%%%%% Stile dei titoli (allineati a sinistra, senza spaziatura verticale extra)
\titleformat{\section}{\normalfont\large\bfseries}{\thesection}{1em}{}
\titleformat{\subsection}{\normalfont\normalsize\bfseries}{\thesubsection}{1em}{}
\titleformat{\subsubsection}{\normalfont\normalsize\bfseries}{\thesubsubsection}{1em}{}
\titleformat{\paragraph}[runin]{\normalfont\normalsize\bfseries}{\theparagraph}{1em}{}[.]
%%%%%%%%%


\usepackage{fancyhdr}         % Intestazione e piè di pagina personalizzati
\pagestyle{fancy}
\fancyhf{}                    % Azzera header e footer predefiniti (evita duplicati)
\fancyhead[LE,RO]{\thepage}           % Numero pagina: sinistra su pari, destra su dispari
\fancyhead[RE]{\nouppercase{\leftmark}}  % Nome capitolo nelle pagine pari
\fancyhead[LO]{\nouppercase{\rightmark}} % Nome sezione nelle pagine dispari
\fancyfoot{}                  % Footer vuoto (il numero di pagina è già nell'header)

% Profondità indice e numerazione sezioni (allineati allo stesso livello)
\setcounter{tocdepth}{2}
\setcounter{secnumdepth}{2}

% =====================================================================
% === 3. MATEMATICA E TEOREMI
% =====================================================================

\usepackage{amsmath, amssymb, amsthm}
\usepackage{mathtools}        % Estensione di amsmath (es. \coloneqq, casi allineati)
\usepackage{bm}               % Grassetto matematico (\bm{v} per vettori)

% --- Stili degli ambienti teorema ---
% "plain"  → testo in corsivo (per Teoremi, Lemmi, Proposizioni)
% "definition" → testo normale (per Definizioni, Esempi)
% "remark" → testo normale, titolo non in grassetto (per Osservazioni)

\theoremstyle{plain}
\newtheorem{theorem}{Teorema}[section]
\newtheorem{lemma}[theorem]{Lemma}
\newtheorem{proposition}[theorem]{Proposizione}
\newtheorem{corollary}[theorem]{Corollario}

\theoremstyle{definition}
\newtheorem{definition}[theorem]{Definizione}
\newtheorem{example}[theorem]{Esempio}

\theoremstyle{remark}
\newtheorem*{osservazione}{Osservazione}

% --- Comando per teoremi in stile libero (uso sconsigliato: non numerati né referenziabili) ---
% Preferire gli ambienti sopra. Lasciato per retrocompatibilità con il documento esistente.
\newcommand{\mytheo}[2]{%
    \noindent \textbf{#1}\par #2 \par \vspace{1.5em}
}

% =====================================================================
% === 4. COLORI
% =====================================================================
% xcolor va caricato prima di colortbl, tcolorbox e altri pacchetti che usano colori.

\usepackage[table, xcdraw]{xcolor}

% --- Palette colori del progetto ---
\definecolor{tableheadbg}{gray}{0.90}
\definecolor{tablerowblue}{rgb}{0.93, 0.95, 1.0}
\definecolor{annotationred}{rgb}{0.8, 0, 0}

% Colori per evidenziazione matrici
\definecolor{matrixhighlight}{RGB}{200, 220, 240}
\definecolor{matrixblue}{RGB}{210, 220, 240}
\definecolor{matrixyellow}{RGB}{250, 250, 210}  % ← unica definizione (era duplicata)

% Colori di sfondo generici
\definecolor{LightMint}{HTML}{E0FFF0}
\definecolor{CoolGray}{HTML}{EAEAEA}

% --- Comandi scorciatoia per testo colorato ---
\newcommand{\bluetext}[1]{\textcolor{blue}{#1}}
\newcommand{\redtext}[1]{\textcolor{annotationred}{#1}}

% =====================================================================
% === 5. TABELLE
% =====================================================================
% Caricare dopo xcolor.

\usepackage{colortbl}         % Celle e righe colorate (dipende da xcolor)
\usepackage{booktabs}         % Tabelle professionali (\toprule, \midrule, \bottomrule)
\usepackage{array}            % Formattazione avanzata colonne (es. >{...}c)
\usepackage{multirow}         % Celle su più righe (\multirow)
\usepackage{nicematrix}       % Matrici avanzate con TikZ integrato

% =====================================================================
% === 6. GRAFICA E FIGURE
% =====================================================================

\usepackage{graphicx}         % Inclusione immagini (\includegraphics)
\usepackage{wrapfig}          % Testo che scorre attorno alle figure
                              % ATTENZIONE: può dare problemi vicino a liste o float

% NOTA: svg richiede Inkscape installato sulla macchina (non funziona su Overleaf di default).
% Per disabilitarlo commentare la riga e usare immagini pre-convertite in PDF/PNG.
\usepackage[inkscapelatex=false]{svg}

% =====================================================================
% === 7. BOX COLORATI (tcolorbox)
% =====================================================================

\usepackage{tcolorbox}
\tcbuselibrary{skins, breakable} % Bordi avanzati e divisione automatica tra pagine

% --- Box "Più in generale" (generalizzazione di un concetto) ---
\newtcolorbox{generalizationbox}{
    colback  = blue!5,
    colframe = blue!50!black,
    title    = {Più in generale},
    fonttitle = \bfseries,
    coltitle = black,
    boxrule  = 0.8pt,
    arc      = 3pt,
    left=6pt, right=6pt, top=4pt, bottom=4pt
}

% --- Box "Esempio" con titolo variabile: \begin{esempio}{Titolo} ---
\newtcolorbox{esempio}[1]{
    enhanced,
    breakable,
    before skip = 12pt,
    after skip  = 12pt,
    boxrule     = 0.5pt,
    left=8pt, right=8pt, top=8pt, bottom=8pt,
    colback     = gray!5,
    colframe    = gray!40,
    sharp corners,
    borderline west = {2pt}{0pt}{blue!70!black},
    title       = \textbf{Esempio: #1},
    coltitle    = blue!70!black,
    attach title to upper,
    after title = {\par\smallskip},
    fontupper   = \small
}

% =====================================================================
% === 8. GRAFICA VETTORIALE (TikZ)
% =====================================================================

\usepackage{tikz}

% --- Librerie TikZ ---
\usetikzlibrary{
    arrows.meta,                  % Frecce avanzate (sintassi moderna: -Stealth)
    positioning,                  % Posizionamento relativo (right=of nodo)
    matrix,                       % Matrici di nodi
    fit,                          % Nodi che racchiudono altri nodi
    tikzmark,                     % Marcatori nel testo per connessioni TikZ
    decorations.pathmorphing,     % Decorazioni di percorso (ondulato, ecc.)
    decorations.pathreplacing,    % Decorazioni con sostituzione (brace, ecc.)
    shapes,                       % Forme generiche
    shapes.callouts,              % Fumetti
    shapes.geometric,             % Forme geometriche (ellisse, diamante, ecc.)
    calc                          % Calcoli su coordinate
}

% --- Stili globali TikZ ---
\tikzset{
    % FIX: usata la sintassi moderna di arrows.meta (Stealth maiuscolo, senza apostrofo)
    edge/.style       = {->, >=Stealth, semithick},
    labelstyle/.style = {auto, font=\small}
}

% =====================================================================
% === 9. ALGORITMI E CODICE SORGENTE
% =====================================================================

\usepackage{algorithm}         % Ambiente float per algoritmi
\usepackage{algpseudocode}     % Comandi per pseudocodice (if, for, while, ecc.)
\usepackage{seqsplit}          % Spezza stringhe molto lunghe (es. sequenze di bit)

\usepackage{listings}          % Inclusione codice sorgente con syntax highlighting
\lstset{
    basicstyle  = \ttfamily\small,   % Font monospazio, corpo leggermente ridotto
    breaklines  = true,              % Manda a capo le righe troppo lunghe
    frame       = single,            % Cornice attorno al codice
    numbers     = left,              % Numeri di riga a sinistra
    numberstyle = \tiny\color{gray}, % Numeri di riga in grigio
    tabsize     = 4,
    showstringspaces = false         % Non mostrare il simbolo ␣ nelle stringhe
}

% =====================================================================
% === 10. UTILITÀ GENERALI
% =====================================================================

\usepackage{comment}           % Blocchi di commento multipagina (\begin{comment})
\usepackage{enumitem}          % Liste personalizzabili (spaziatura, simboli, ecc.)
\usepackage{etoolbox}
\AtBeginEnvironment{algorithmic}{\setlist{}}  % isola enumitem dagli ambienti algoritmo
\usepackage{placeins}          % \FloatBarrier per bloccare i float in una sezione
\usepackage[normalem]{ulem}    % Sottolineature avanzate (\uline, \sout, ecc.)
                               % [normalem] evita che ridefinisca \emph
% NOTA: soul (\hl per evidenziare) e ulem ridefiniscono entrambi \underline internamente.
% In caso di conflitti, commentare uno dei due e usare solo l'altro.
\usepackage{soul}              % Evidenziazione (\hl), spaziatura lettere (\so)

% =====================================================================
% === 11. NOTE REDAZIONALI (todonotes)
% =====================================================================


% =====================================================================
% === 11. RIFERIMENTI E LINK (ULTIMI DA CARICARE)
% =====================================================================
% hyperref va caricato quasi sempre per ultimo.
% cleveref DEVE essere caricato dopo hyperref.

\usepackage[
    colorlinks = true,
    linkcolor  = blue,
    urlcolor   = blue,
    citecolor  = blue
]{hyperref}


% Usare \cref{label} al posto di \ref{label}.
\usepackage[italian]{cleveref}
